\documentclass[11pt]{article}
\usepackage{amsmath, amssymb, geometry, mathtools, booktabs, longtable}
\usepackage{hyperref}
\geometry{margin=1in}

\title{RG Checkerboard Flow with In-Block Mixing for 2D $\phi^4$: Code-Level Description and CLI Reference}
\author{Kostas Orginos}
\date{\today}

\begin{document}
\maketitle

\section{Purpose}

This note documents the in-block checkerboard RG flow implemented in
\texttt{jaxqft/models/phi4\_rg\_checkerboard\_inblock\_flow.py}, together with
its trainer
\texttt{scripts/phi4/train\_rg\_checkerboard\_inblock\_flow.py}
and checker
\texttt{scripts/phi4/check\_rg\_checkerboard\_inblock\_flow.py}.

This path is a new evolution branch. It leaves
\texttt{jaxqft/models/phi4\_rg\_checkerboard\_flow.py},
\texttt{jaxqft/models/phi4\_rg\_cond\_flow.py}, and
\texttt{jaxqft/models/phi4\_mg.py} unchanged.

\section{High-Level Design}

The model combines four ingredients:
\begin{itemize}
\item recursive $2\times 2$ RG splitting into one coarse mode and three local
fluctuation modes,
\item red/black checkerboard updates acting on block fluctuations,
\item a shifted second blocking obtained by rolling the fine lattice by
$(1,1)$ before repeating the checkerboard pass,
\item masked in-block RealNVP-style mixing among the three fluctuation
components inside each red or black color pass.
\end{itemize}

Relative to the earlier checkerboard branch, the essential new feature is that
the three fluctuation components
\[
\eta=(\eta_1,\eta_2,\eta_3)\in\mathbb{R}^3
\]
inside one block are no longer updated all at once with no internal coupling.
Instead, each red or black pass is itself decomposed into a sequence of masked
subupdates that mix the three components while still conditioning only on
frozen variables.

\section{Lattice, Prior, and RG Split}

Let the fine lattice be
\[
\Lambda_0,\qquad |\Lambda_0| = L\times L,\qquad L=2^K.
\]
The batched field is
\[
\phi \in \mathbb{R}^{B\times L\times L}.
\]

The latent prior is sitewise standard normal:
\begin{equation}
p_Z(z)=\prod_{x\in\Lambda_0}\frac{1}{\sqrt{2\pi}}
\exp\!\left(-\frac{z_x^2}{2}\right).
\end{equation}

At RG level $k$, the current lattice has size
\[
L_k = L/2^k.
\]
Each $2\times 2$ block is packed into
\[
v=(v_{00},v_{01},v_{10},v_{11})\in\mathbb{R}^4.
\]

\subsection{Average Blocking}

For \texttt{rg\_type="average"}, the code uses the orthonormal Haar-like basis
\begin{equation}
H=
\begin{pmatrix}
\tfrac12 & \tfrac12 & \tfrac12 & \tfrac12 \\
\tfrac12 & -\tfrac12 & \tfrac12 & -\tfrac12 \\
\tfrac12 & \tfrac12 & -\tfrac12 & -\tfrac12 \\
\tfrac12 & -\tfrac12 & -\tfrac12 & \tfrac12
\end{pmatrix}
\end{equation}
and defines
\begin{equation}
(c,\eta_1,\eta_2,\eta_3)=H\,v.
\end{equation}
Thus
\[
\eta=(\eta_1,\eta_2,\eta_3)
\]
spans the local fluctuation subspace.

The arithmetic block average is
\[
\bar v = \frac14\sum_{i\in\mathrm{block}} v_i = \frac12\,c.
\]
This arithmetic average is the coarse scalar used in conditioning.

\subsection{Select Blocking}

For \texttt{rg\_type="select"}, the code uses
\begin{equation}
c=v_{00},\qquad \eta=(v_{01},v_{10},v_{11}).
\end{equation}

\subsection{Inverse Reconstruction}

At each block the inverse reconstruction is
\[
v = R(c,\eta).
\]
For average blocking, this is simply
\[
v = H^\top(c,\eta).
\]

\section{Recursive Structure}

The recursion stops at the final $2\times 2$ lattice, where the remaining
$4$-component block is mapped by an unconditional RealNVP.

For each non-terminal level:
\begin{enumerate}
\item apply a full-field level bijection $T_k$ on the current lattice,
\item split the result into a next coarse lattice and stored fine modes,
\item recurse on the coarse lattice,
\item reconstruct the fine lattice by merging the recursively transformed
coarse field with the stored fine modes.
\end{enumerate}

So, as in the previous checkerboard branch, the checkerboard dynamics is a
\emph{full-field} map at each scale, not a post hoc correction applied only to
previously extracted $\eta$ variables.

\section{Checkerboard Plus In-Block Level Map}

Fix one RG level and suppress the level index. Let
\[
S(x) = (c,\eta)
\]
denote the block split, where $c$ lives on the block lattice and
\[
\eta_b\in\mathbb{R}^3
\]
is attached to each block $b$.

Let the block lattice be decomposed into red and black checkerboards:
\[
\Lambda_{\mathrm{blk}} = R \cup B,\qquad R\cap B = \varnothing.
\]
The corresponding masks are denoted by $M_R$ and $M_B$.

\subsection{In-Block Masks}

Inside each block, the code uses the familiar $3$-component coupling masks
\[
m^{(1)}=(1,1,0),\qquad
m^{(2)}=(0,1,1),\qquad
m^{(3)}=(1,0,1),
\]
or the first \texttt{n\_inner\_couplings} of this sequence if that integer is
smaller than three. Their complements are
\[
\bar m^{(i)} = \mathbf 1 - m^{(i)}.
\]

For mask $m^{(i)}$:
\begin{itemize}
\item the components selected by $m^{(i)}$ are visible to the conditioner,
\item the complementary components selected by $\bar m^{(i)}$ are updated.
\end{itemize}

\subsection{One Red Color Pass}

Fix one in-block mask $m^{(i)}$. In a red pass:
\begin{itemize}
\item all black blocks are frozen and visible,
\item on red blocks, only the components selected by $m^{(i)}$ are visible,
\item on red blocks, only the complementary components selected by
$\bar m^{(i)}$ are updated.
\end{itemize}

Define the visible mask
\begin{equation}
V_R^{(i)} = M_B + M_R\,m^{(i)},
\end{equation}
and the update mask
\begin{equation}
U_R^{(i)} = M_R\,\bar m^{(i)}.
\end{equation}

The conditioner input is
\begin{equation}
u_R^{(i)} = \bigl[V_R^{(i)}\eta,\ \tilde c\bigr],
\end{equation}
where
\[
\tilde c =
\begin{cases}
\tfrac12 c,& \text{average blocking,}\\
c,& \text{select blocking.}
\end{cases}
\]

Then the affine subupdate is
\begin{align}
s_R^{(i)},t_R^{(i)} &= \mathcal C_R^{(i)}(u_R^{(i)}), \\
\eta &\mapsto (1-U_R^{(i)})\eta
+ U_R^{(i)}\bigl(\eta\odot e^{s_R^{(i)}} + t_R^{(i)}\bigr),
\end{align}
with the code-level stabilization
\[
s_R^{(i)} \leftarrow \texttt{log\_scale\_clip}\cdot \tanh(s_R^{(i)}).
\]

The full red color pass composes these subupdates in order
\[
i=1,\dots,N_{\mathrm{inner}},
\qquad
N_{\mathrm{inner}}=\texttt{n\_inner\_couplings}.
\]

\subsection{One Black Color Pass}

The black pass is analogous, with red and black exchanged:
\begin{equation}
V_B^{(i)} = M_R + M_B\,m^{(i)},
\qquad
U_B^{(i)} = M_B\,\bar m^{(i)}.
\end{equation}
The black pass uses the already-updated red variables from the preceding red
pass as part of its frozen conditioning field.

\subsection{One Checkerboard Pass}

One checkerboard pass on the current blocking is therefore
\begin{equation}
P = P_B \circ P_R,
\end{equation}
where each of $P_R$ and $P_B$ is itself a composition of
\texttt{n\_inner\_couplings} masked in-block affine subupdates.

\subsection{Shifted Blocking and One Full Cycle}

After one checkerboard pass, the field is merged back to the fine lattice,
rolled by $(-1,-1)$ with periodic wrapping, updated again on the shifted
blocking, and rolled back. One cycle is
\begin{equation}
T = \rho^{-1}\circ P_{\mathrm{shift}}\circ \rho \circ P_{\mathrm{orig}}.
\end{equation}

The code repeats this cycle
\[
\texttt{n\_cycles}
\]
times per non-terminal RG level.

\section{Why the Map Remains Invertible}

The critical condition is unchanged:
\begin{quote}
every affine subupdate must depend only on variables that remain fixed during
that subupdate.
\end{quote}

This is enforced at two levels.

\subsection{Across Blocks}

When updating red blocks, the full black checkerboard is visible but frozen.
When updating black blocks, the full red checkerboard is visible but frozen.

\subsection{Inside One Block}

For a fixed in-block mask $m^{(i)}$, the updated components are the ones
selected by $\bar m^{(i)}$, and these are masked out of the conditioner input
everywhere. On the same-color blocks, the conditioner sees only the held-fixed
components selected by $m^{(i)}$.

\subsection{Attention Radius}

Therefore increasing \texttt{attn\_radius} does \emph{not} break
invertibility, provided the masking is applied before attention, as in the
current code. A larger radius only exposes a larger neighborhood of frozen
variables.

So:
\begin{itemize}
\item larger \texttt{attn\_radius} increases cost,
\item larger \texttt{attn\_radius} increases the frozen conditioning context,
\item larger \texttt{attn\_radius} does not violate the triangular coupling
structure in the current implementation.
\end{itemize}

\section{Local Conditioner}

The code supports
\[
\texttt{conditioner}\in\{\texttt{transformer},\texttt{mlp}\}.
\]

\subsection{MLP Conditioner}

For \texttt{conditioner="mlp"}, each masked subupdate uses a pointwise MLP
applied independently at each block site:
\[
(s,t)=\mathrm{MLP}(u_b).
\]
There is no spatial attention.

\subsection{Transformer Conditioner}

For \texttt{conditioner="transformer"}, each masked subupdate uses a local
transformer on the block lattice. The token at block $b$ is
\[
u_b\in\mathbb{R}^4,
\]
consisting of the masked visible fluctuation vector in $\mathbb{R}^3$ plus the
coarse scalar.

For \texttt{attn\_radius} $= r$, the attended neighborhood is
\[
\mathcal N_r(b)=\{b+\delta:\delta_y,\delta_x\in\{-r,\dots,r\}\}.
\]
Hence:
\begin{itemize}
\item \texttt{attn\_radius = 0} means on-site conditioning only,
\item \texttt{attn\_radius = 1} means a $3\times 3$ neighborhood on the block
lattice,
\item in general the receptive field contains $(2r+1)^2$ block tokens.
\end{itemize}

\subsection{$Z_2$ Symmetry}

The code supports
\[
\texttt{parity}\in\{\texttt{none},\texttt{sym}\}.
\]

For \texttt{parity="sym"}, the raw conditioner outputs are projected to
\begin{align}
s(u) &= \frac12\bigl(\hat s(u) + \hat s(-u)\bigr), \\
t(u) &= \frac12\bigl(\hat t(u) - \hat t(-u)\bigr),
\end{align}
so that
\[
s(-u)=s(u),\qquad t(-u)=-t(u).
\]
This preserves the old-map-style symmetry
\[
z\mapsto -z,\qquad \phi\mapsto -\phi.
\]

\section{Terminal $2\times 2$ Flow}

At the coarsest $2\times 2$ lattice, the remaining four real variables are
packed into a vector in $\mathbb{R}^4$ and mapped by an unconditional RealNVP
\[
g_{\mathrm{term}}:\mathbb{R}^4\to\mathbb{R}^4.
\]
This terminal flow uses the same parity option as the in-block checkerboard
conditioners.

\section{Exact Log Density}

The inverse map returns
\[
(z,\log|\det J^{-1}|).
\]
Hence
\begin{equation}
\log p_\theta(x)
= \log p_Z(z)
+ \log|\det J_{\mathrm{term}}^{-1}|
+ \sum_k \log|\det J_k^{-1}|.
\end{equation}

Each masked affine subupdate contributes a diagonal Jacobian on the updated
subset of variables only. The total log-Jacobian at one level is the sum over:
\begin{itemize}
\item both checkerboard colors,
\item both blockings in each cycle,
\item all \texttt{n\_inner\_couplings} masked in-block subupdates.
\end{itemize}

\section{Training Objective}

Training uses the standard flow objective
\begin{equation}
\mathcal L(\theta)=
\mathbb E_{z\sim p_Z}
\bigl[
\log p_\theta(g_\theta(z)) + S(g_\theta(z))
\bigr].
\end{equation}

For reweighting quality, the practically relevant validation quantity is the
centered width of
\[
\Delta S = \log p_\theta(\phi) + S(\phi),
\]
namely
\[
\sigma\bigl(\Delta S - \langle\Delta S\rangle\bigr),
\]
because the reweighting factors are
\[
w = e^{-\Delta S}.
\]

\section{CLI Reference: Training Script}

This section documents
\texttt{scripts/phi4/train\_rg\_checkerboard\_inblock\_flow.py}.

\begin{longtable}{p{0.24\linewidth}p{0.70\linewidth}}
\toprule
Argument & Meaning \\
\midrule
\endhead
\texttt{--resume} & Resume training from an existing checkpoint. \\
\texttt{--save} & Output checkpoint path. \\
\texttt{--save-every} & Save an intermediate checkpoint every $N$ epochs. Zero means save only at the end. \\
\texttt{--validate} & Run the final validation routine after training. \\
\texttt{--validate-batch} & Batch size used inside one validation draw. \\
\texttt{--validate-super} & Number of validation draws concatenated before computing $\Delta S$ and ESS summaries. \\
\texttt{--L} & Lattice linear size; the lattice is $L\times L$. Must be a power of two. \\
\texttt{--lam} & Quartic coupling $\lambda$. \\
\texttt{--mass} & Quadratic mass parameter $m^2$. \\
\texttt{--width} & Conditioner width. For the transformer, this is the model dimension. For the MLP, this is the hidden width. \\
\texttt{--n-cycles} & Number of original-blocking plus shifted-blocking checkerboard cycles per non-terminal RG level. \\
\texttt{--n-inner-couplings} & Number of masked in-block affine subupdates inside each red or black color pass. \\
\texttt{--conditioner} & Conditioner family: \texttt{transformer} or \texttt{mlp}. \\
\texttt{--transformer-layers} & Number of transformer blocks in each masked subupdate network. Used only for \texttt{conditioner="transformer"}. \\
\texttt{--n-heads} & Number of attention heads. Used only for \texttt{conditioner="transformer"}. \\
\texttt{--attn-radius} & Local attention radius on the block lattice. Used only for \texttt{conditioner="transformer"}. \\
\texttt{--mlp-dim} & Hidden width of the transformer MLP blocks. Used only for \texttt{conditioner="transformer"}. Zero means $2\times$\texttt{width}. \\
\texttt{--terminal-n-layers} & Number of RealNVP layers in the terminal $2\times 2$ flow. \\
\texttt{--terminal-width} & Hidden width of the terminal RealNVP MLPs. Zero means reuse \texttt{width}. \\
\texttt{--output-init-scale} & Small nonzero output initialization scale used to start near the identity while still allowing gradients to flow. \\
\texttt{--parity} & Symmetry mode. \texttt{sym} enforces the $\phi\to-\phi$, $z\to-z$ symmetry; \texttt{none} does not. \\
\texttt{--rg-type} & RG blocking rule: \texttt{average} or \texttt{select}. \\
\texttt{--log-scale-clip} & Tanh clip for affine log-scales. \\
\texttt{--epochs} & Total target epoch count. On resume this is the final epoch number, not the number of extra epochs. \\
\texttt{--batch} & Training batch size. \\
\texttt{--lr} & Adam learning rate. \\
\texttt{--seed} & Random seed. \\
\bottomrule
\end{longtable}

\section{CLI Reference: Checker Script}

This section documents
\texttt{scripts/phi4/check\_rg\_checkerboard\_inblock\_flow.py}.

\begin{longtable}{p{0.24\linewidth}p{0.70\linewidth}}
\toprule
Argument & Meaning \\
\midrule
\endhead
\texttt{--shape} & Lattice shape used for invertibility and timing tests. \\
\texttt{--jac-shape} & Small lattice shape used for the dense autodiff Jacobian check. \\
\texttt{--batch-size} & Batch size used in invertibility and timing tests. \\
\texttt{--width}, \texttt{--n-cycles}, \texttt{--n-inner-couplings}, \texttt{--conditioner}, \texttt{--transformer-layers}, \texttt{--n-heads}, \texttt{--attn-radius}, \texttt{--mlp-dim}, \texttt{--terminal-n-layers}, \texttt{--terminal-width}, \texttt{--output-init-scale}, \texttt{--parity}, \texttt{--rg-type}, \texttt{--log-scale-clip} & Same architectural meanings as in the training script. \\
\texttt{--lam}, \texttt{--mass} & Theory parameters used in the timing loss-gradient benchmark. \\
\texttt{--tests} & Comma-separated subset of \texttt{invertibility}, \texttt{jacobian}, \texttt{timing}. \\
\texttt{--invert-tol} & Absolute tolerance for the forward-then-inverse reconstruction test. \\
\texttt{--jac-tol} & Tolerance for the autodiff-vs-analytic logdet comparison. \\
\texttt{--jac-perturb-scale} & Small perturbation applied to initialized weights before the Jacobian test, so the logdet is not trivially zero. \\
\texttt{--timing-warmup} & Warmup iterations for each timed kernel. \\
\texttt{--timing-iters} & Number of timed iterations per kernel. \\
\texttt{--json-out} & Optional JSON output path for the checker summary. \\
\texttt{--selfcheck-fail} & Exit with nonzero status if any selected test fails. \\
\bottomrule
\end{longtable}

\section{Practical Interpretation}

This architecture is a direct attempt to fix the main structural limitation of
the earlier checkerboard branch:
\begin{quote}
the old checkerboard map coupled each block to its neighbors, but it did not
mix the three fluctuation components inside the block during one color pass.
\end{quote}

The new path addresses this by nesting masked in-block couplings inside each
color pass. So it now mixes:
\begin{itemize}
\item across neighboring blocks through the checkerboard and shifted-blocking
mechanism,
\item within one block through the sequence of masked in-block subupdates.
\end{itemize}

This makes the map more expressive, but it also increases cost. In practice,
the first tuning knob should usually be \texttt{conditioner="mlp"} if one wants
to isolate the effect of the in-block coupling structure before paying the full
cost of the transformer.

\end{document}
